\originalsection{受迫摆为什么需要三维状态}
以无量纲时间$t=\sqrt{g/\ell}\,t_{\rm physical}$, 受迫阻尼摆写成$\dot\theta=\omega$, $\dot\omega=-\omega/q-\sin\theta+a\cos(\omega_Dt)$, $q>0$. 若只记$(\theta,\omega)$, 方程显含时间, 相平面同一点在不同驱动相位受到不同力. 引入$\varphi=\omega_Dt\pmod{2\pi}$后成为自治系统
\[
\dot\theta=\omega,\qquad
\dot\omega=-\frac\omega q-\sin\theta+a\cos\varphi,\qquad
\dot\varphi=\omega_D.
\]
真实状态是$(\theta,\omega,\varphi)$. 平面投影交叉不违反微分方程解的唯一性, 因为驱动相位可能不同. 散度为$-1/q$, 三维相体积收缩; 但收缩本身不能区分稳定周期轨道与混沌运动. 小角近似把$\sin\theta$换成$\theta$, 仍保留周期驱动, 即可有稳定周期响应; 并不是驱动项必须是某种“非线性力”才有孤立周期轨道.

取驱动周期$T_D=2\pi/\omega_D$, 每次在相同$\varphi$采样, 得到映射$P:(\theta_n,\omega_n)\mapsto(\theta_{n+1},\omega_{n+1})$, 称频闪Poincaré映射. 长时间后一个点对应与驱动同周期的响应, 两点交替对应二倍周期, $k$点对应$kT_D$周期. 若把角度模$2\pi$, 还须记录连续角的变化, 区分原位摆动与每周期转过整数圈的旋转. 闭合曲线状点集常提示准周期, 散布而复杂的点集可提示混沌, 二者都应结合重复采样、数值收敛和扰动增长检验; 一张点图不能独自证明结论. 第17章作业中的RLC实验给出了实际积分的频闪数据, 并单列参数、暂态、采样长度与误差检查.

本章采用的动力学意义是: 确定性系统在有关有界不变集合上出现非周期运动及初值敏感性. 局部不稳定且向无穷远逃逸的$\dot x=x$虽有误差指数增长, 不作为此处的混沌实例. 若初误差$\delta_0$在相关时间段按$\delta(t)\simeq\delta_0e^{\lambda t}$增长, 容许误差为$\Delta$, 可预测时间约为$\lambda^{-1}\log(\Delta/\delta_0)$. 这个估计依赖线性误差区间, 饱和后不能继续无限外推.

\originalsection{滞回与物理阈值}
带稳定高阶项的亚临界叉形模型$\dot x=\mu x+x^3-x^5$说明, 局部三次项之外的结构会改变全局结果. 非零固定点令$y=x^2$, 得$y^2-y-\mu=0$, 故
\[
y_\pm=\frac{1\pm\sqrt{1+4\mu}}2.
\]
在$-1/4<\mu<0$时两个$y$均为正, 每个给一对对称固定点. 非零点处导数为$2y(1-2y)$, 所以外侧$y_+$稳定、内侧$y_-$不稳定; 原点导数$\mu<0$, 也稳定. 缓慢增大参数时, 原点在$\mu=0$失稳后可跳到外侧稳定支; 缓慢减小时, 外侧支一直保留到$\mu=-1/4$的鞍结. 因而两个转变点不同, 称滞回. 恰好位于原点的理想确定性解在参数变化后仍可形式上留在原点, 实际跳转需要微扰; 不能把稳定性改变和自动离开固定点混为一谈.

原课的激光模型为$\dot n=GnN-Kn$, $\dot N=-GnN-fN+p$, 参数均正. 假设激发粒子数$N$比光子数$n$快得多地弛豫, 先用$N\simeq p/(Gn+f)$, 才得到一维近似
\[
\dot n=n\left(\frac{Gp}{Gn+f}-K\right)
=\left(\frac{Gp}f-K\right)n-\frac{G^2p}{f^2}n^2+O(n^3).
\]
物理区域为$n\ge0$. 泵强度阈值为$p_c=Kf/G$; 阈值以上非零根的精确表达为$n_*=p/K-f/G$. 小振幅截断模型的根是$r/b$, 其中$r=Gp/f-K$, $b=G^2p/f^2$, 不能漏掉$b$写成$n_*=r$. 原稿第159页这里须按模型系数修正. 这一推导是快慢变量近似的阈值模型, 没有证明原始激光系统的全部瞬态行为.

\originalsection{极限环、中心与环形陷获区}
极限环是孤立周期轨道: 邻近没有一族任意接近的周期轨道. 无阻尼谐振子每个能量都给一个椭圆, 属于中心的周期轨道族, 不称极限环. 极限环可以稳定、失稳, 也可以只从一侧吸引. 把超临界模型改为$\dot r=\mu r+r^3-r^5$, $\dot\theta=\omega_0\ne0$, 径向方程的正根正是上一节的$\sqrt{y_\pm}$; 由相同导数可判定外圈吸引、内圈排斥. 在$-1/4<\mu<0$时, 稳定原点与稳定外环共存, 中间不稳定环分开两个吸引域. 这给亚临界Hopf和滞回的一种具体模型, 不是所有二维系统的固定公式.

Van der Pol方程$\ddot x-\mu(1-x^2)\dot x+x=0$, $\mu>0$, 写成$\dot x=y$, $\dot y=\mu(1-x^2)y-x$. 能量函数$E=(x^2+y^2)/2$满足$\dot E=\mu(1-x^2)y^2$. 小$|x|$时输入能量, 大$|x|$时耗散能量, 解释为何振幅可被调节到有限值. 仅凭这个符号变化尚未证明极限环存在与唯一; 原课在此用数值相图展示, 本册不把图形补成未经证明的唯一性定理.

要真正调用Poincaré--Bendixson, 可构造一个可核算的陷获区. 原课例子
\[
\dot x=x-y-x(x^2+5y^2),\qquad
\dot y=x+y-y(x^2+y^2)
\]
给$r=\sqrt{x^2+y^2}$的导数$\dot r=r[1-r^2(1+\sin^22\theta)]$. 在$r=0.65<1/\sqrt2$上径向导数严格正, 在$r=1.05>1$上严格负, 故环域$0.65\le r\le1.05$正向不变. 它没有固定点: 若$y=0$, 第二式迫使$x=0$; 若$y\ne0$, 令$d=1-r^2$, 第二式给$x=-dy$, 第一式化为$d^2-4dy^2+1=0$. $d<0$时左端正; $0\le d\le1$时$y^2\le1-d$, 故左端至少$5d^2-4d+1\ge1/5>0$. 因而唯一固定点是环外的原点. 环内任一轨道有非空紧极限集, 不含固定点, 定理保证其中有周期轨道. 这个论证没有额外证明周期轨道唯一, 也不能把有鞍点或连接轨道的任意平面系统都归成“固定点或极限环”两类.

恒转矩摆$\dot\theta=\omega$, $\dot\omega=\tau-\sin\theta-\omega/q$也可用圆柱相空间分析. 对$\tau>1$, 取$m=q(\tau-1)/2$、$M=q(\tau+1)+1$. 在$\omega=m$上$\dot\omega>0$, 在$\omega=M$上$\dot\omega<0$, 形成$\theta\pmod{2\pi}$上的紧环带. 且$\dot\theta>0$, 没有固定点, 由对应环带的平面嵌入可得旋转周期轨道存在. 能量$E=\omega^2/2-\cos\theta$满足$\dot E=\tau\omega-\omega^2/q$, 一圈后$2\pi\tau=q^{-1}\int\omega^2\dd t$: 输入功与耗散相平衡. 对$0<\tau<1$固定点可与旋转运动共存, 是否共存还依赖$q$; 原课示意稳定区边界, 未给全参数精确阈值公式.

\originalsection{伸展、收缩与Lyapunov谱}
自治系统$\dot z=F(z)$的微小扰动满足$\dot\eta=DF(z(t))\eta$. 令基本矩阵$Y(0)=I$, $\dot Y=DF(z(t))Y$, 对长时间伸展倍率的奇异值取对数除以$t$, 在极限存在时得到Lyapunov指数$\lambda_1\ge\cdots\ge\lambda_n$. 单个方向的倍率依赖方向及不断变化的主轴; 数值计算常需要反复正交化, 避免所有向量都挤到最快方向而失去其他指数.

Jacobi行列式满足$\dd\log|\det Y|/\dd t=\operatorname{tr}DF=\nabla\cdot F$. 因奇异值的乘积为$|\det Y|$, 有
\[
\sum_i\lambda_i
=\lim_{t\to\infty}\frac1t\int_0^t\nabla\cdot F(z(s))\dd s.
\]
这使某方向伸展与总体积收缩可以同时成立. 受迫阻尼摆的三维和为$-1/q$, Lorenz系统为$-(\sigma+1+\beta)$; Hamilton流的和为零. 非平衡自治轨道沿流方向的微扰只是时间平移, 在通常有界、非退化条件下对应零指数, 不应把稳定周期轨道的相位方向误判成吸引固定点的所有方向衰减. 有限时指数、极限指数、集合是否吸引及是否有分形结构仍是不同问题.

\originalsection{Lorenz的有界性与奇异吸引子}
Lorenz方程的两个非零固定点及其局部稳定分析见作业9第3题. 例如标准参数的原点不稳定, 不等于已经证明出现混沌. 要理解耗散轨道为何可能留在有限区域, 可补一个直接估计. 记$C=\rho+\sigma$, $V=x^2+y^2+(z-C)^2$, 对原方程求导, 三次交叉项抵消, 得
\[
\frac12\dot V=-\sigma x^2-y^2-\beta z(z-C).
\]
以$w=z-C$, 用$w^2+Cw\ge w^2/2-C^2/2$, 得$\dot V\le-2mV+\beta C^2$, $m=\min(\sigma,1,\beta/2)>0$. 乘积分因子$e^{2mt}$后积分,
\[
V(t)\le V(0)e^{-2mt}
 +\frac{\beta C^2}{2m}(1-e^{-2mt}).
\]
于是所有有限初值都保持有界, 大椭球最终陷获轨道. 这是中文补充的有界性证明, 没有推导出该区域内一定有混沌. 标准参数的复杂双翼轨道、初值敏感性和负散度共同提示奇异吸引子结构, 其严格证明超出此课的推导范围. 也不能由一次参数下的模拟推断所有超过局部Hopf阈值的$\rho$都混沌.

\originalsection{分形构造与维数的限度}
Cantor构造第$n$步总区间长度为$(2/3)^n\to0$, 因而极限集长度为零; 每步仍留下$2^n$个自相似部分, 得盒维数$\log2/\log3$. Koch曲线每段替成四条长度为原来$1/3$的折线, 第$n$步总折线长度为$(4/3)^n$, 发散; 以边长$3^{-n}$的小盒覆盖, 数量尺度为$4^n$, 总覆盖面积尺度为$(4/9)^n\to0$, 维数为$\log4/\log3$. 这些是明确几何构造, 不含动力学敏感依赖; 分形不能作为混沌的定义.

\begin{center}
\begin{tikzpicture}[x=7cm,y=.48cm]
\node[anchor=east,font=\small] at (-.04,0) {$n=0$};
\draw[blue!65!black,line width=2pt] (0.0000000000,0)--(1.0000000000,0);
\node[anchor=east,font=\small] at (-.04,-1) {$n=1$};
\draw[blue!65!black,line width=2pt] (0.0000000000,-1)--(0.3333333333,-1);
\draw[blue!65!black,line width=2pt] (0.6666666667,-1)--(1.0000000000,-1);
\node[anchor=east,font=\small] at (-.04,-2) {$n=2$};
\draw[blue!65!black,line width=2pt] (0.0000000000,-2)--(0.1111111111,-2);
\draw[blue!65!black,line width=2pt] (0.2222222222,-2)--(0.3333333333,-2);
\draw[blue!65!black,line width=2pt] (0.6666666667,-2)--(0.7777777778,-2);
\draw[blue!65!black,line width=2pt] (0.8888888889,-2)--(1.0000000000,-2);
\node[anchor=east,font=\small] at (-.04,-3) {$n=3$};
\draw[blue!65!black,line width=2pt] (0.0000000000,-3)--(0.0370370370,-3);
\draw[blue!65!black,line width=2pt] (0.0740740741,-3)--(0.1111111111,-3);
\draw[blue!65!black,line width=2pt] (0.2222222222,-3)--(0.2592592593,-3);
\draw[blue!65!black,line width=2pt] (0.2962962963,-3)--(0.3333333333,-3);
\draw[blue!65!black,line width=2pt] (0.6666666667,-3)--(0.7037037037,-3);
\draw[blue!65!black,line width=2pt] (0.7407407407,-3)--(0.7777777778,-3);
\draw[blue!65!black,line width=2pt] (0.8888888889,-3)--(0.9259259259,-3);
\draw[blue!65!black,line width=2pt] (0.9629629630,-3)--(1.0000000000,-3);
\node[font=\small] at (.5,.8) {Cantor构造};
\end{tikzpicture}
\par\smallskip
\begin{tikzpicture}[x=7cm,y=7cm]
\draw[blue!65!black,thin] (0.0000000000,0.0000000000)--(1.0000000000,0.0000000000);
\node[anchor=east,font=\small] at (-.04,-0.0000000000) {$n=0$};
\draw[blue!65!black,thin] (0.0000000000,-0.3000000000)--(0.3333333333,-0.3000000000)--(0.5000000000,-0.0113248654)--(0.6666666667,-0.3000000000)--(1.0000000000,-0.3000000000);
\node[anchor=east,font=\small] at (-.04,-0.3000000000) {$n=1$};
\draw[blue!65!black,thin] (0.0000000000,-0.6000000000)--(0.1111111111,-0.6000000000)--(0.1666666667,-0.5037749551)--(0.2222222222,-0.6000000000)--(0.3333333333,-0.6000000000)--(0.3888888889,-0.5037749551)--(0.3333333333,-0.4075499103)--(0.4444444444,-0.4075499103)--(0.5000000000,-0.3113248654)--(0.5555555556,-0.4075499103)--(0.6666666667,-0.4075499103)--(0.6111111111,-0.5037749551)--(0.6666666667,-0.6000000000)--(0.7777777778,-0.6000000000)--(0.8333333333,-0.5037749551)--(0.8888888889,-0.6000000000)--(1.0000000000,-0.6000000000);
\node[anchor=east,font=\small] at (-.04,-0.6000000000) {$n=2$};
\draw[blue!65!black,thin] (0.0000000000,-0.9000000000)--(0.0370370370,-0.9000000000)--(0.0555555556,-0.8679249850)--(0.0740740741,-0.9000000000)--(0.1111111111,-0.9000000000)--(0.1296296296,-0.8679249850)--(0.1111111111,-0.8358499701)--(0.1481481481,-0.8358499701)--(0.1666666667,-0.8037749551)--(0.1851851852,-0.8358499701)--(0.2222222222,-0.8358499701)--(0.2037037037,-0.8679249850)--(0.2222222222,-0.9000000000)--(0.2592592593,-0.9000000000)--(0.2777777778,-0.8679249850)--(0.2962962963,-0.9000000000)--(0.3333333333,-0.9000000000)--(0.3518518519,-0.8679249850)--(0.3333333333,-0.8358499701)--(0.3703703704,-0.8358499701)--(0.3888888889,-0.8037749551)--(0.3703703704,-0.7716999402)--(0.3333333333,-0.7716999402)--(0.3518518519,-0.7396249252)--(0.3333333333,-0.7075499103)--(0.3703703704,-0.7075499103)--(0.3888888889,-0.6754748953)--(0.4074074074,-0.7075499103)--(0.4444444444,-0.7075499103)--(0.4629629630,-0.6754748953)--(0.4444444444,-0.6433998804)--(0.4814814815,-0.6433998804)--(0.5000000000,-0.6113248654)--(0.5185185185,-0.6433998804)--(0.5555555556,-0.6433998804)--(0.5370370370,-0.6754748953)--(0.5555555556,-0.7075499103)--(0.5925925926,-0.7075499103)--(0.6111111111,-0.6754748953)--(0.6296296296,-0.7075499103)--(0.6666666667,-0.7075499103)--(0.6481481481,-0.7396249252)--(0.6666666667,-0.7716999402)--(0.6296296296,-0.7716999402)--(0.6111111111,-0.8037749551)--(0.6296296296,-0.8358499701)--(0.6666666667,-0.8358499701)--(0.6481481481,-0.8679249850)--(0.6666666667,-0.9000000000)--(0.7037037037,-0.9000000000)--(0.7222222222,-0.8679249850)--(0.7407407407,-0.9000000000)--(0.7777777778,-0.9000000000)--(0.7962962963,-0.8679249850)--(0.7777777778,-0.8358499701)--(0.8148148148,-0.8358499701)--(0.8333333333,-0.8037749551)--(0.8518518519,-0.8358499701)--(0.8888888889,-0.8358499701)--(0.8703703704,-0.8679249850)--(0.8888888889,-0.9000000000)--(0.9259259259,-0.9000000000)--(0.9444444444,-0.8679249850)--(0.9629629630,-0.9000000000)--(1.0000000000,-0.9000000000);
\node[anchor=east,font=\small] at (-.04,-0.9000000000) {$n=3$};
\node[font=\small] at (.5,.23) {Koch曲线构造};
\end{tikzpicture}
\end{center}

有限级图形只有有限线段, 不是真的非整数维集合. 维数结论来自任意尺度的构造规律及覆盖估计, 图示用于帮助理解迭代. 盒维数与Hausdorff维数的定义不同, 在这两个满足分离条件的自相似例子中数值一致; 不能据此把任意吸引子的两种维数混为一谈.

原课还用伸展和压缩估计吸引子维数. 常用Kaplan--Yorke预测为$D_{\rm KY}=j+(\lambda_1+\cdots+\lambda_j)/|\lambda_{j+1}|$, 其中$j$为前$j$项和非负而前$j+1$项和负的最大指标. 若三维谱为$\lambda_1>0$, $\lambda_2=0$, $\lambda_3<0$且总和负, 则$D_{\rm KY}=2+\lambda_1/|\lambda_3|$. 二维示意模型才会出现$1+\lambda_1/|\lambda_2|$. 这是基于局部伸展率的维数预测, 没有普遍证明它等于每个系统的盒维数或Hausdorff维数.

Logistic的倍周期分岔点$r_n$若趋向$r_\infty$, 间距比$(r_n-r_{n-1})/(r_{n+1}-r_n)$在相应光滑单峰映射类中趋于约$4.669$, 是原课介绍的Feigenbaum标度. 几何缩小的间距可帮助数值外推累积参数, 但本册不证明该普适常数; 超过累积点仍有周期窗口, “达到第一次混沌阈值后一直混沌”不是结论.

\originalsection{湍流的惯性区间与量纲分析}
三维高Re统计稳态湍流的大尺度驱动输入能量, 小尺度黏性耗散能量. 在两者之间的惯性区间, 设每单位质量、每单位时间传递的能量为$\varepsilon$, 量纲为$L^2T^{-3}$. 假设该区间近似局部各向同性, 且尺度$\ell$的速度差只由$\varepsilon,\ell$控制, 量纲给$v_\ell\sim(\varepsilon\ell)^{1/3}$. 这一条件说明涡周转时间$\ell/v_\ell$内每单位质量能量$v_\ell^2$的传递率尺度为$v_\ell^3/\ell$.

黏性在最小尺度$\eta_K$上不再可忽略, 令尺度Re为$v_{\eta_K}\eta_K/\nu\sim1$, 与前式联立得
\[
\eta_K\sim\left(\frac{\nu^3}{\varepsilon}\right)^{1/4},\qquad
v_{\eta_K}\sim(\nu\varepsilon)^{1/4},\qquad
t_{\eta_K}\sim\left(\frac\nu\varepsilon\right)^{1/2}.
\]
这里$\eta_K$为长度, 切勿与上一章的动力黏度$\eta$混同. 若能谱定义为每单位质量动能$\int E(k)\dd k$, 则$[E]=L^3T^{-2}$, 只用$k$和$\varepsilon$构造它, 给Kolmogorov惯性区间形式$E(k)=C_K\varepsilon^{2/3}k^{-5/3}$. 量纲分析确定指数, 不确定无量纲常数$C_K$, 也没有由此解出全部速度场、边界效应或间歇性修正. 范围须有$L^{-1}\ll k\ll\eta_K^{-1}$; 二维湍流的传递结构不同, 不能不加条件套用此三维图景.
